\documentclass{article}%
\usepackage[T1]{fontenc}%
\usepackage[utf8]{inputenc}%
\usepackage{lmodern}%
\usepackage{textcomp}%
\usepackage{lastpage}%
%
\usepackage{amsmath}%
\usepackage{amssymb}%
\usepackage{amsfonts}%
\usepackage{graphicx}%
\usepackage{hyperref}%
\usepackage{tikz}%
\usepackage{pgfplots}%
\usetikzlibrary{math}%
\title{IME 2027 - Física (one-shot, sem judge)\thanks{This exam was solved and automatically generated with gpt-resolve: \url{https://github.com/lgabs/gpt-resolve}}}%
\author{claude{-}opus{-}5{-}5}%
\date{October 30, 2026}%
%
\begin{document}%
\normalsize%
\maketitle%
\section*{Questão 1}

\section*{Solução}

\textbf{Hipóteses.} Sem resistência do ar, depois de sair do lançador o foguete é um projétil sob $g$ constante (lançamento oblíquo). Origem na plataforma, eixo $x$ horizontal na direção do alvo e eixo $y$ vertical para cima. O alvo fica em $(2000;\,1000)$ m.

\textbf{Conferência do projeto.} Com $v_0=200$ m/s e $\theta=45^\circ$, temos $v_{0x}=v_{0y}=100\sqrt2$ m/s. Então
\[
H=\frac{v_{0y}^2}{2g}=\frac{20000}{20}=1000\ \text{m},\qquad
x_H=\frac{v_{0x}v_{0y}}{g}=\frac{20000}{10}=2000\ \text{m},
\]
e o vértice coincide de fato com o centro do alvo. Os dados são coerentes.

\textbf{a) Velocidade de lançamento na falha.} Eliminando o tempo em $x=v\cos\theta\,t$ e $y=v\sin\theta\,t-\tfrac12gt^2$, obtém-se a equação da trajetória:
\[
y=x\tan\theta-\frac{g\,x^2}{2v^2\cos^2\theta}\;\xrightarrow{\ \theta=45^\circ\ }\; y=x-\frac{g\,x^2}{v^2}.
\]
O ângulo não mudou. Na vertical do alvo, $x=2000$ m e $y=1000+360=1360$ m:
\[
1360=2000-\frac{10\cdot(2000)^2}{v^2}\;\Rightarrow\; v^2=\frac{4\cdot 10^{7}}{640}=62500\ \text{m}^2/\text{s}^2
\;\Rightarrow\; \boxed{v=250\ \text{m/s}}.
\]
\emph{Verificação do movimento ascendente:} o vértice da nova trajetória está em $x_H'=\dfrac{v^2\sin 2\theta}{2g}=\dfrac{62500}{20}=3125\ \text{m}>2000\ \text{m}$. Logo, em $x=2000$ m o foguete ainda sobe. Explicitamente, $v_y=v\sin\theta-g\,\dfrac{x}{v\cos\theta}=125\sqrt2-80\sqrt2=45\sqrt2\approx 63{,}6\ \text{m/s}>0$, o que confirma o enunciado.

\textbf{b) Temperatura atingida.} Pela observação, $T\propto v^2$:
\[
\frac{T}{T_0}=\left(\frac{v}{v_0}\right)^2\;\Rightarrow\;
T=2560\cdot\left(\frac{250}{200}\right)^2=2560\cdot 1{,}5625\;\Rightarrow\;\boxed{T=4000\ \text{K}}.
\]
Portanto, $\Delta T=4000-2560=1440$ K.

\textbf{c) Trabalho no aquecimento.} No aquecimento isobárico, o calor recebido pelos $m=2{,}0$ kg de gás é
\[
Q=m\,c_p\,\Delta T=2{,}0\cdot 1{,}0\cdot 1440=2880\ \text{kJ}.
\]
Pela 1.ª Lei da Termodinâmica, $Q=\Delta U+W$, com $\Delta U=2016$ kJ:
\[
W=Q-\Delta U=2880-2016\;\Rightarrow\;\boxed{W=864\ \text{kJ}}.
\]
O sinal positivo indica trabalho realizado \emph{pelos} gases, que se expandem ao serem aquecidos a pressão constante.

\emph{Consistência:} $c_v=\dfrac{\Delta U}{m\,\Delta T}=\dfrac{2016}{2\cdot1440}=0{,}70$ kJ/(kg$\cdot$K). Daí $\gamma=c_p/c_v\approx1{,}43$, valor típico de gás diatômico. Além disso, $W=p\,\Delta V=m\,(c_p-c_v)\,\Delta T=2\cdot0{,}30\cdot1440=864$ kJ, que confere com o resultado acima.

ANSWER: a) $v=250$ m/s (a $45^\circ$); b) $T=4000$ K; c) $W=864$ kJ (realizado pelos gases).%
\newpage%
\section*{Questão 2}

\section*{Solução}

\textbf{Hipóteses e geometria.} Trata-se de uma treliça ideal: as barras têm peso desprezível, são articuladas nas extremidades e a carga atua só nos nós. Cada barra é, portanto, um elemento de duas forças, e o esforço normal $N$ atua ao longo do eixo da barra. Convenção: $N>0$ indica tração e $N<0$, compressão. Com origem em $A$: $A(0,0)$, $B(3,0)$, $C(-3,4)$, $D(0,4)$, $E(3,4)$, $F(0,8)$ (em m). As barras inclinadas $CF$, $EF$ e $AE$ medem $\sqrt{3^2+4^2}=5$ m, com cossenos diretores $3/5$ (horizontal) e $4/5$ (vertical).

Contagem: $b+r=8+4=12=2n=2\cdot 6$. O problema é isostático e todas as incógnitas saem do equilíbrio.

\textbf{a) Reações.} A barra $BE$ é a única ligada a $B$ e é vertical. Como ela só transmite força ao longo do próprio eixo, a reação em $B$ é vertical: $B_x=0$. Para a estrutura inteira (momentos anti-horários positivos):
\[
\sum F_x=0:\ A_x+B_x=0\ \Rightarrow\ A_x=0,
\]
\[
\sum M_A=0:\ (-3)(-P)+3\,B_y=0\ \Rightarrow\ B_y=-P=-100\ \text{kN},
\]
\[
\sum F_y=0:\ A_y+B_y-P=0\ \Rightarrow\ A_y=2P=200\ \text{kN}.
\]
Logo $\vec R_A$ tem $200$ kN, é vertical e aponta para cima. Já $\vec R_B$ tem $100$ kN, é vertical e aponta para baixo: o apoio $B$ precisa ancorar (puxar) a estrutura.

\textbf{b) Método dos nós.} Em cada nó, a força de uma barra sobre ele vale $N$ vezes o vetor unitário que aponta do nó para a outra extremidade da barra.

\emph{Nó C:}
\[
\tfrac45 N_{CF}-P=0\Rightarrow N_{CF}=125\ \text{kN};\qquad N_{CD}+\tfrac35N_{CF}=0\Rightarrow N_{CD}=-75\ \text{kN}.
\]
\emph{Nó F:}
\[
-\tfrac35N_{CF}+\tfrac35N_{EF}=0\Rightarrow N_{EF}=125\ \text{kN};\qquad
-\tfrac45N_{CF}-N_{DF}-\tfrac45N_{EF}=0\Rightarrow N_{DF}=-200\ \text{kN}.
\]
\emph{Nó D} (só há barras horizontais e verticais):
\[
N_{DE}=N_{CD}=-75\ \text{kN};\qquad N_{AD}=N_{DF}=-200\ \text{kN}.
\]
\emph{Nó E:}
\[
\textstyle\sum F_x:\ -N_{DE}-\tfrac35N_{EF}-\tfrac35N_{AE}=0\Rightarrow 75-75-\tfrac35N_{AE}=0\Rightarrow N_{AE}=0,
\]
\[
\textstyle\sum F_y:\ \tfrac45N_{EF}-\tfrac45N_{AE}-N_{BE}=0\Rightarrow N_{BE}=100\ \text{kN}.
\]
\emph{Verificação nos apoios.} Em $B$, a barra $BE$ tracionada com $100$ kN puxa o nó para cima, e o apoio reage com $100$ kN para baixo. Em $A$, a barra $AD$ comprimida com $200$ kN empurra o nó para baixo, e como $N_{AE}=0$ o apoio reage com $A_y=200$ kN e $A_x=0$. Os resultados coincidem com o item (a).

\begin{center}
\begin{tabular}{c|c|c}
Barra & $|N|$ (kN) & Esforço\\ \hline
$CD$ & $75$ & compressão\\
$DE$ & $75$ & compressão\\
$CF$ & $125$ & tração\\
$EF$ & $125$ & tração\\
$DF$ & $200$ & compressão\\
$AD$ & $200$ & compressão\\
$BE$ & $100$ & tração\\
$AE$ & $0$ & nulo
\end{tabular}
\end{center}

\textbf{c) Função da barra AE.} Sem $AE$, teríamos $b+r=7+4=11<12$: a estrutura seria hipostática. A parte superior $CDEF$, rígida por ser triangulada, ficaria apoiada apenas em duas barras verticais paralelas, $AD$ e $BE$, articuladas no solo. O painel retangular $ADEB$ formaria então um mecanismo de paralelogramo, e a parte superior poderia oscilar lateralmente. Além disso, nesse deslocamento a carga desce, de modo que o equilíbrio seria instável e a estrutura colapsaria.

A diagonal $AE$ triangula esse painel e torna a estrutura rígida e isostática: é uma barra de contraventamento (travamento lateral). Com a carga puramente vertical deste problema ela não trabalha ($N_{AE}=0$). Ainda assim, é indispensável para a estabilidade e para resistir a qualquer ação horizontal, como vento, componente horizontal da carga ou perturbações dinâmicas.

\medskip
ANSWER: a) $\vec R_A$: $200$ kN, vertical, para cima; $\vec R_B$: $100$ kN, vertical, para baixo (componentes horizontais nulas). b) $CD=DE=75$ kN (compressão); $CF=EF=125$ kN (tração); $DF=AD=200$ kN (compressão); $BE=100$ kN (tração); $AE=0$. c) $AE$ é a diagonal de contraventamento: triangula o painel $ADEB$ e impede o mecanismo de deslocamento lateral. Sem esforço sob esta carga vertical, ela garante a rigidez e a estabilidade da estrutura.%
\newpage%
\section*{Questão 3}

\section*{Solução}

\textbf{Hipóteses.} Tomamos $g = GM/R^2 = 12\ \text{m/s}^2$. A haste é vertical e a argola não tem atrito, então a força que ela aplica na barra é vertical. Toda a irradiância que chega à face superior é absorvida e vira calor sensível. Desprezamos as perdas por emissão e o trabalho de dilatação.

\medskip
\textbf{Massa e peso.} O volume é $V = 2\cdot 0{,}25\cdot 0{,}05 = 0{,}025\ \text{m}^3$. Logo
\[
m = \rho V = 2000\cdot 0{,}025 = 50\ \text{kg},\qquad P = mg = 50\cdot 12 = 600\ \text{N}.
\]
A massa não muda com o aquecimento, e portanto o peso também não.

\medskip
\textbf{a) Reação no apoio ao amanhecer.} Nenhuma força externa tem componente horizontal: a tração é vertical e a argola não tem atrito. Assim, a reação $N$ do apoio articulado também é vertical. A distribuição de massa é homogênea, então o centro de massa fica em $L/2$, e a haste atua em $x = L$. Tomando momentos em relação ao apoio:
\[
T_0 L = P\,\frac{L}{2}\ \Rightarrow\ T_0 = \frac{P}{2} = 300\ \text{N},
\qquad
N + T_0 = P\ \Rightarrow\ N = 300\ \text{N}\ (\text{vertical, para cima}).
\]

\medskip
\textbf{b) Variação de temperatura até o rompimento.} A extremidade esquerda continua presa ao apoio, e a haste continua em $x = L$. A barra desliza para a direita pela argola e passa a ter comprimento $L' = L(1+\alpha\,\Delta\theta)$. Como a massa continua homogênea, o centro de massa vai para $x = L'/2$. Os braços de alavanca são horizontais, por isso a dilatação na espessura não interfere. Momentos em relação ao apoio:
\[
T L = P\,\frac{L(1+\alpha\Delta\theta)}{2}\ \Rightarrow\ T = T_0\,(1+\alpha\,\Delta\theta).
\]
A haste se rompe quando $T = 1{,}001\,T_0 = 300{,}3\ \text{N}$. Então
\[
\alpha\,\Delta\theta = 10^{-3}\ \Rightarrow\ \Delta\theta = \frac{10^{-3}}{10^{-5}} = 100\ ^\circ\text{C}.
\]
Nesse instante a barra mede $2{,}002$ m e sobram $2$ mm à direita da haste.

\medskip
\textbf{c) Horário do rompimento.} Com $c = 1{,}0\ \text{J/(g}\,^\circ\text{C)} = 1000\ \text{J/(kg}\,^\circ\text{C)}$, o calor necessário é
\[
Q = m\,c\,\Delta\theta = 50\cdot 1000\cdot 100 = 5\cdot 10^{6}\ \text{J}.
\]
A área que absorve é a face superior: $A = 2\cdot 0{,}25 = 0{,}5\ \text{m}^2$.

O trecho de 25\% a 50\% do dia dura $\Delta t_1 = 0{,}25\cdot 8\cdot 10^4 = 2\cdot 10^4$ s. Nesse trecho, sendo $\tau$ o tempo medido desde o amanhecer,
\[
I(\tau) = \frac{4000}{2\cdot 10^4}\,\tau = 0{,}2\,\tau\ \ \text{(W/m}^2\text{, com } \tau \text{ em s)}.
\]
Até o pico ($\tau = 2\cdot 10^4$ s), a energia absorvida é a área do triângulo vezes $A$:
\[
E_{\max} = \frac{1}{2}\cdot 2\cdot 10^4\cdot 4000\cdot 0{,}5 = 2\cdot 10^7\ \text{J}.
\]
Como $E_{\max} > Q$, o rompimento acontece ainda na rampa de subida. A energia absorvida até $\tau$ é
\[
E(\tau) = A\int_0^{\tau} 0{,}2\,\tau'\,d\tau' = 0{,}5\cdot 0{,}1\,\tau^2 = 0{,}05\,\tau^2 .
\]
Igualando a $Q$:
\[
0{,}05\,\tau^2 = 5\cdot 10^6\ \Rightarrow\ \tau = 10^4\ \text{s} = \frac{10^4}{8\cdot 10^4} = 12{,}5\%\ \text{do dia}.
\]
O horário do rompimento é, portanto,
\[
t = 25\% + 12{,}5\% = 37{,}5\%\ \text{do dia}.
\]

\textit{Verificação.} Nesse instante a irradiância vale $0{,}2\cdot 10^4 = 2000\ \text{W/m}^2$, metade do pico. Isso é coerente com o fato de $E \propto \tau^2$: um quarto da energia da rampa corresponde à metade do seu tempo, já que $5\cdot 10^6 = \frac{1}{4}\cdot 2\cdot 10^7$.

O aumento da área superior com a dilatação, $A = A_0(1+\alpha\Delta\theta)^2$, muda a área em no máximo $0{,}2\%$. Isso antecipa o rompimento em cerca de $5$ s, ou seja, $0{,}006\%$ do dia, o que é desprezível.

\medskip
ANSWER: a) $N = 300$ N, vertical e para cima; b) $\Delta\theta = 100\ ^\circ$C; c) $37{,}5\%$ do dia.%
\newpage%
\section*{Questão 4}

\section*{Solução}

\textbf{Hipóteses.} Fusível, chave e fios são ideais (resistência desprezível). Os receptores seguem o modelo dado: f.c.e.m. constante em série com $R_{int}$, ou seja, $U = E' + R_{int}\,i$. Chamamos de $V = V_A - V_B$ a tensão comum aos três ramos, que estão todos ligados entre os nós $A$ e $B$.

\subsection*{a) Correntes logo após o fechamento de S}

Equações de ramo, com $I$ saindo do gerador e $i_1, i_2$ entrando nos receptores:
\[
I=\frac{1000-V}{2},\qquad i_1=\frac{V-850}{1},\qquad i_2=\frac{V-900}{2}.
\]
Lei dos nós em $A$: $I=i_1+i_2$. Isso equivale ao teorema de Millman:
\[
V=\frac{\dfrac{1000}{2}+\dfrac{850}{1}+\dfrac{900}{2}}{\dfrac12+\dfrac11+\dfrac12}=\frac{500+850+450}{2}=900\ \text{V}.
\]
Logo
\[
I=\frac{1000-900}{2}=50\ \text{A},\qquad i_1=900-850=50\ \text{A},\qquad i_2=\frac{900-900}{2}=0 .
\]
O valor $i_2=0$ é coerente: a tensão disponível ($900$ V) apenas iguala a f.c.e.m. de Rec$_2$, que, portanto, não recebe corrente.

Verificação de potência:
\begin{itemize}
\item potência fornecida pelo gerador: $1000\cdot 50=50000$ W;
\item potência dissipada na resistência interna do gerador: $2\cdot 50^2=5000$ W;
\item potência absorvida por Rec$_1$: $850\cdot 50+1\cdot 50^2=45000$ W.
\end{itemize}
Como $5000+45000=50000$ W, o balanço fecha.

\subsection*{b) Rompimento do fusível}

Toda a corrente do gerador passa pelo fusível:
\[
\frac{I}{I_{FUS}}=\frac{50}{20}=2{,}5 .
\]
Na Figura 2, a assíntota vertical em $I/I_{FUS}\approx 1{,}1$ indica que só há rompimento para razões maiores que cerca de $1{,}1$. Como $2{,}5>1{,}1$, o fusível \textbf{vai romper}. As correntes são constantes, porque o modelo é resistivo com f.c.e.m. constante, então o tempo de rompimento é lido diretamente na curva.

Na região de interesse, a curva é um segmento de reta que vai de aproximadamente $(1{,}1;\ 3{,}8\ \text{s})$ até $(3{,}3;\ 0{,}1\ \text{s})$. Interpolando:
\[
t\approx 3{,}8-\frac{3{,}8-0{,}1}{3{,}3-1{,}1}\,(2{,}5-1{,}1)\approx 3{,}8-1{,}68\cdot 1{,}4\approx 1{,}5\ \text{s}.
\]
Portanto, o fusível rompe cerca de $1{,}5$ s após o fechamento de S. Esse valor vem de leitura gráfica, com incerteza de aproximadamente $\pm 0{,}1$ s. A partir desse instante, o gerador fica desligado dos receptores.

\subsection*{c) Curto-circuito entre A e B}

Com o curto ideal, $V_{AB}=0$. Os receptores ficam em paralelo com o curto e não alteram a corrente do ramo do gerador: as correntes que eles eventualmente fornecessem circulariam pelo próprio curto, sem passar por F. O gerador, então, opera em curto-circuito:
\[
I_{cc}=\frac{\varepsilon}{r}=\frac{1000}{2}=500\ \text{A}\quad\Rightarrow\quad \frac{I_{cc}}{I_{FUS}}=\frac{500}{20}=25 .
\]
Na Figura 2, para $I/I_{FUS}\gtrsim 3{,}3$ a curva é horizontal no valor mínimo $t=10^{-1}$ s. Esse patamar representa o tempo mínimo de atuação do fusível e vale também para a razão $25$, que está fora da escala do gráfico. Logo
\[
t_{rompimento}=0{,}1\ \text{s}=100\ \text{ms}>50\ \text{ms}.
\]
O gerador ficaria em curto durante $100$ ms, o dobro do tempo máximo que suporta. Assim, o fusível \textbf{não} consegue proteger a fonte. Seria necessário um dispositivo com tempo de atuação de no máximo $50$ ms para $I=500$ A.

ANSWER: a) $V_{AB}=900$ V; $i_{\text{Rec}_1}=50$ A, $i_{\text{Rec}_2}=0$ (corrente no gerador e no fusível $I=50$ A). b) Sim: $I/I_{FUS}=2{,}5>1{,}1$, e o fusível rompe em $t\approx 1{,}5$ s (leitura gráfica). c) $I_{cc}=500$ A, $I/I_{FUS}=25$, e o fusível rompe em $0{,}1$ s $=100$ ms $>50$ ms; logo não protege o gerador.%
\newpage%
\section*{Questão 5}

\section*{Solução}

\textbf{Hipóteses.} As duas máquinas operam em regime permanente (ciclo), então trabalhamos com taxas (potências). Valem a 1ª lei aplicada ao ciclo ($\dot W=\dot Q_{\text{recebido}}-\sum\dot Q_{\text{rejeitado}}$) e a 2ª lei na forma de Clausius para vários reservatórios: $\sum \dot Q_i/T_i\le 0$, com $\dot Q_i$ positivo quando o calor entra no ciclo. O ``máximo admissível'' entre dois reservatórios é o rendimento de Carnot.

\medskip
\textbf{a) Temperatura da fonte fria de A.}
Potência útil: $\dot W_A=80\cdot\tfrac34=60$ kW. Rendimento real:
\[
\eta_A=\frac{\dot W_A}{\dot Q_H}=\frac{60}{150}=0{,}40 .
\]
Como $\eta_A=0{,}8\,\eta_{\text{Carnot}}$, temos $\eta_{\text{Carnot}}=0{,}40/0{,}8=0{,}50$. Daí
\[
1-\frac{T}{900}=0{,}50\;\Longrightarrow\; T=450\ \text{K}.
\]

\textbf{b) Calor rejeitado por A.} Pela 1ª lei aplicada ao ciclo:
\[
\dot Q=\dot Q_H-\dot W_A=150-60=90\ \text{kW}.
\]
Verificação pela 2ª lei: $-\tfrac{150}{900}+\tfrac{90}{450}=\tfrac{1}{30}\ \text{kW/K}>0$, ou seja, há geração de entropia, como se espera de uma máquina real (irreversível).

\textbf{c) Análise da proposta B.}
Dados de B: $\dot Q_H=150$ kW a $900$ K e $\dot Q/2=45$ kW rejeitados a $T=450$ K; $\dot Q_3$ é rejeitado a $300$ K.

\emph{Demandas anunciadas.} Cada residência consome em média $240/720=\tfrac13$ kW, logo
\[
\dot W_{\text{elet}}=350\cdot\tfrac13=\tfrac{350}{3}\approx 116{,}7\ \text{kW}.
\]
Para a água, com $\dot m=\tfrac{50}{60}\ \text{kg/s}=\tfrac56$ kg/s:
\[
\dot Q_3=\dot m\,c\,\Delta T=\tfrac56\cdot 4000\cdot 10\ \text{W}=\tfrac{100}{3}\ \text{kW}\approx 33{,}3\ \text{kW}.
\]

\emph{1ª lei.} O ciclo exigiria
\[
\dot W+\dot Q/2+\dot Q_3=116{,}7+45+33{,}3=195\ \text{kW}>150\ \text{kW},
\]
o que viola a conservação de energia. Note que $116{,}7+33{,}3=150$ kW: a proposta simplesmente ignora os $45$ kW rejeitados para a fonte a $450$ K. Mesmo sem aquecer a água, o balanço daria $\dot Q_3=150-45-116{,}7=-11{,}7$ kW. Isso significa que o ciclo teria de \emph{retirar} calor da fonte a $300$ K, o que contradiz o próprio funcionamento proposto. Nesse caso a 2ª lei também falharia: $-\tfrac{150}{900}+\tfrac{45}{450}-\tfrac{11{,}7}{300}\approx-0{,}106\ \text{kW/K}<0$.

\emph{Limite pela 2ª lei (Clausius).}
\[
-\frac{150}{900}+\frac{45}{450}+\frac{\dot Q_3}{300}\le 0\;\Longrightarrow\;\dot Q_3\ge 300\left(\tfrac16-\tfrac1{10}\right)=20\ \text{kW}.
\]
Combinando com a 1ª lei, $\dot W=105-\dot Q_3\le 85$ kW. Esse é o limite reversível, que atenderia no máximo $85/(1/3)=255$ residências. Portanto, as $350$ residências são impossíveis mesmo para uma máquina ideal.

\emph{O que B pode de fato entregar.} Aquecer a água é possível, pois $\dot Q_3=33{,}3\ \text{kW}\ge 20$ kW. Nessas condições,
\[
\dot W=150-45-33{,}3\approx 71{,}7\ \text{kW},
\]
com geração de entropia $-\tfrac16+\tfrac1{10}+\tfrac{33{,}3}{300}\approx 0{,}044\ \text{kW/K}>0$, que é fisicamente aceitável. Isso atende $71{,}7\cdot 3=215$ residências e ainda supera a máquina A ($60$ kW, isto é, $180$ residências).

\emph{Conclusão.} As alegações de B não são realistas. Atender $350$ residências e aquecer $50$ L/min ao mesmo tempo exigiria $195$ kW de saídas para apenas $150$ kW de entrada, violando a 1ª lei. Só a parte elétrica ($116{,}7$ kW) já supera o máximo de $85$ kW permitido pela 2ª lei. O que B consegue de forma realista é aquecer a água e fornecer cerca de $71{,}7$ kW, o que atende $215$ residências.

\medskip
ANSWER: a) $T=450$ K; b) $\dot Q=90$ kW; c) as alegações de B não são realistas. Seriam necessários $116{,}7+33{,}3+45=195$ kW $>150$ kW (viola a 1ª lei), e a 2ª lei limita $\dot W\le 85$ kW ($\le 255$ residências). Aquecendo a água ($33{,}3$ kW), B fornece cerca de $71{,}7$ kW, o que atende $215$ residências.%
\newpage%
\section*{Questão 6}

\section*{Solução}

\textbf{Hipóteses.} O som se propaga no ar ($v_{som}=340$ m/s), com o observador O em repouso e no plano da trajetória. A frequência percebida depende da componente da velocidade da fonte, \emph{no instante da emissão}, ao longo da reta fonte--observador:
\[
f=f_0\,\frac{v_{som}}{v_{som}-v\cos\theta},
\]
em que $\theta$ é o ângulo entre $\vec v$ e o vetor que vai da fonte até O. Adotamos eixos na vista superior com origem em A, $x$ apontando para O e $y$ apontando para B. Assim, $A=(0,0)$, o centro é $C=(0,R)$, $B=(0,2R)$ e $O=(d,0)$.

\textbf{a) Velocidade escalar.} Como o movimento é anti-horário, em A (ponto mais baixo) a velocidade é tangente à circunferência e aponta no sentido $+x$, isto é, exatamente na direção de O ($\theta=0$, já que AO é tangente em A):
\[
\frac{5f_0}{4}=f_0\,\frac{340}{340-v}\;\Rightarrow\;1700-5v=1360\;\Rightarrow\; v=68\ \text{m/s}.
\]
Como $v/v_{som}=0{,}2<1$, a fonte é subsônica e a fórmula vale. O resultado é coerente com $f_A>f_0$, pois a fonte se aproxima de O.

\textbf{b) Raio.} Em B a velocidade é $\vec v=(-v,0)$, e o vetor que vai de B até O é $(d,-2R)$. A projeção de $\vec v$ nessa direção é $-v\cos\theta$, com $\cos\theta=\dfrac{d}{\sqrt{d^2+4R^2}}$. Logo a fonte se afasta de O:
\[
\frac{10f_0}{11}=f_0\,\frac{340}{340+v\cos\theta}\;\Rightarrow\;340+v\cos\theta=374\;\Rightarrow\;\cos\theta=\frac{34}{68}=\frac12 .
\]
Portanto
\[
\frac{d}{\sqrt{d^2+4R^2}}=\frac12\;\Rightarrow\;4R^2=3d^2\;\Rightarrow\;R=\frac{\sqrt3}{2}\,d=5\sqrt3\ \text{m}\approx 8{,}66\ \text{m}.
\]

\textbf{c) Tração.} O fio forma com a vertical um ângulo $\alpha$ tal que $\sin\alpha=\dfrac{R}{L}=\dfrac{\sqrt3}{2}$, ou seja, $\alpha=60^\circ$ e $\cos\alpha=\dfrac12$. A profundidade da piscina é $L\cos\alpha=5$ m.

Sobre o objeto atuam o peso $mg$ (para baixo), o empuxo $E$ (vertical, para cima) e a tração $T$ ao longo do fio (para baixo e para o eixo). Não há resistência da água. Como o MCU é horizontal, a única componente radial é $T\sin\alpha$, que faz o papel de força centrípeta:
\[
T\sin\alpha=\frac{mv^2}{R}\;\Rightarrow\;
T=\frac{\sqrt3\cdot 68^2}{5\sqrt3\cdot\frac{\sqrt3}{2}}=\frac{8}{5}\cdot\frac{34^2}{\sqrt3}=\frac{8}{5}\cdot 667{,}4\approx 1067{,}8\ \text{N}.
\]

\textbf{d) Volume submerso.} Não há aceleração vertical, logo a resultante vertical é nula:
\[
E=mg+T\cos\alpha=10\sqrt3+\frac{1067{,}84}{2}=17{,}32+533{,}92=551{,}24\ \text{N}.
\]
Pelo princípio de Arquimedes, $E=\rho\, V_{sub}\, g$:
\[
V_{sub}=\frac{551{,}24}{1000\cdot 10}\approx 5{,}51\times10^{-2}\ \text{m}^3\ (\approx 55{,}1\ \text{L}).
\]

\textbf{Verificação.} As unidades conferem: $T$ e $E$ resultam em N e $V_{sub}$ em m$^3$. O empuxo é muito maior que o peso ($17{,}3$ N) porque precisa equilibrar também a componente vertical da grande tração exigida pela força centrípeta.

ANSWER: a) $v=68$ m/s; b) $R=5\sqrt3\approx 8{,}66$ m; c) $T\approx 1067{,}8$ N; d) $V_{sub}\approx 5{,}51\times10^{-2}$ m$^3$ ($\approx 55{,}1$ L).%
\newpage%
\section*{Questão 7}

\section*{Solução}

\textbf{Geometria e dados.} No dodecágono regular, o ângulo externo vale $360^\circ/12=30^\circ$. Como AB é horizontal, o lado BC desce com inclinação $\theta=30^\circ$. O lado mede $\overline{BC}=d_0+L_0=1{,}15$ m. Usaremos
\[
K_e q_1 q_2 = 9\cdot 10^{9}\cdot 20\cdot 10^{-6}\cdot 5\cdot 10^{-6}=0{,}9\ \text{N}\,\text{m}^2,\qquad mg\,\text{sen}\,30^\circ = 0{,}4\cdot 10\cdot 0{,}5 = 2\ \text{N}.
\]

\textbf{a)} Pela observação do enunciado, P$_2$ fica praticamente fixa. Sobre P$_1$ atuam o peso e a força de Coulomb, ambos conservativos, e a normal, que não realiza trabalho; não há atrito. Em B, a repulsão inicial vale $0{,}9/1^2=0{,}9$ N $<2$ N, logo P$_1$ desce. A aproximação máxima ocorre quando $v_1=0$. Pela conservação da energia entre B e esse ponto (distância $d$ até P$_2$), com descida de $(d_0-d)$ ao longo de BC:
\[
mg\,(d_0-d)\,\text{sen}\,30^\circ = K_e q_1 q_2\left(\frac{1}{d}-\frac{1}{d_0}\right)
\;\Rightarrow\; 2(1-d)=\frac{0{,}9\,(1-d)}{d}.
\]
Descartando $d=1$ m (posição inicial), obtemos $d=\dfrac{K_e q_1 q_2}{mg\,\text{sen}\,30^\circ}=0{,}45$ m.

\textbf{b)} Como a mola é muito rígida, seu período próprio é desprezível frente ao tempo do movimento de P$_1$. Logo P$_2$ acompanha a situação de forma quase estática, e na aproximação máxima está em equilíbrio. A pequena oscilação devida a a mola começar relaxada é desprezada, e por isso se trata de uma estimativa. Ao longo de BC, no sentido de C, com compressão $x$:
\[
kx = mg\,\text{sen}\,30^\circ + \frac{K_e q_1 q_2}{d^2} = 2+\frac{0{,}9}{0{,}2025}=2+\frac{40}{9}=\frac{58}{9}\ \text{N}.
\]
A mola armazena $E_{el}=\tfrac12 kx^2$. P$_2$ desce $x$ ao longo de BC, ou seja, perde a altura $x\,\text{sen}\,30^\circ$, e assim $\Delta U_g=-mg\,x\,\text{sen}\,30^\circ$. Então
\[
\frac{E_{el}}{|\Delta U_g|}=\frac{\tfrac12 kx^2}{mg\,x\,\text{sen}\,30^\circ}=\frac{kx}{2\,mg\,\text{sen}\,30^\circ}=\frac{58/9}{4}=\frac{29}{18}\approx 1{,}6 .
\]
Como $\Delta U_g<0$, a razão com sinal vale $E_{el}/\Delta U_g\approx -1{,}6$. A razão não depende de $k$. Quando $k\to\infty$, as duas energias tendem a zero, o que é coerente com desprezar o deslocamento de P$_2$ no item a).

\textbf{c)} Após o giro de $30^\circ$ no sentido anti-horário, BC fica horizontal. As partículas estão em repouso e a mola foi retirada. As posições são $\overline{BP_1}=d_0-d=0{,}55$ m, $\overline{P_1P_2}=0{,}45$ m e $\overline{P_2C}=L_0=0{,}15$ m (a compressão $x$ é desprezível).

Peso e normal são perpendiculares a BC, e a força elétrica é interna ao par. Portanto a quantidade de movimento ao longo de BC se conserva: $0=mv_1+mv_2$. Como as massas são iguais, as velocidades têm o mesmo módulo e sentidos opostos, e as partículas percorrem distâncias iguais $s$. P$_2$ chega a C quando $s=0{,}15$ m, enquanto P$_1$ precisaria de $0{,}55$ m para chegar a B. Assim, o primeiro vértice atingido é C, com separação $d'=0{,}45+2\cdot 0{,}15=0{,}75$ m.

Pela conservação da energia (a energia potencial gravitacional não varia):
\[
K_e q_1 q_2\left(\frac{1}{d}-\frac{1}{d'}\right)=2\cdot\frac{m v^2}{2}
\;\Rightarrow\; 0{,}9\left(\frac{20}{9}-\frac{4}{3}\right)=0{,}8\ \text{J}=0{,}4\,v^2
\;\Rightarrow\; v=\sqrt{2}\ \text{m/s}\approx 1{,}4\ \text{m/s}.
\]
A velocidade de P$_2$ é horizontal, no sentido de B para C. Nesse instante, P$_1$ tem velocidade de mesmo módulo, no sentido de C para B.

ANSWER: a) $d=0{,}45$ m; b) $E_{el}/|\Delta U_g|=29/18\approx 1{,}6$ (com sinal, $E_{el}/\Delta U_g\approx -1{,}6$, pois $\Delta U_g<0$); c) $v_{P_2}=\sqrt{2}\approx 1{,}4$ m/s, ao atingir o vértice C.%
\newpage%
\section*{Questão 8}

\section*{Solução}

\textbf{Hipóteses e modelo.} A luz incide normalmente na rede e as franjas são estreitas, então valem as condições de máximo principal de Fraunhofer. O feixe é estreito comparado com as distâncias envolvidas, de modo que a rede funciona como uma origem pontual dos feixes difratados, sobre o eixo que passa pelo centro do anteparo.

\textbf{Passo 1 (ordens possíveis).} O passo da rede é
\[
d=\frac{1\ \text{mm}}{400}=2{,}5\ \mu\text{m}, \qquad \frac{d}{\lambda}=\frac{2{,}5\ \mu\text{m}}{0{,}625\ \mu\text{m}}=4 .
\]
Máximos principais: $d\sin\theta_n=n\lambda \Rightarrow \sin\theta_n=\dfrac{n}{4}$. A ordem $n=4$ daria $\theta=90^\circ$, rasante ao plano da rede, e nunca atinge um anteparo paralelo a ela. Logo existem só as ordens $n=0,\pm1,\pm2,\pm3$, ou seja, no máximo $7$ franjas, o que é coerente com as $7$ observadas antes de $t_1$. Com $\cos\theta_n=\sqrt{1-n^2/16}$:
\[
\tan\theta_1=\frac{1}{\sqrt{15}},\qquad \tan\theta_2=\frac{1}{\sqrt{3}},\qquad \tan\theta_3=\frac{3}{\sqrt{7}} .
\]

\textbf{Passo 2 (condição de visibilidade).} Seja $D$ a distância da rede ao anteparo. A franja de ordem $\pm n$ cai a uma altura $y_n=D\tan\theta_n$ do eixo. Ela é vista enquanto $y_n\le L/2=0{,}6$ m, isto é, enquanto
\[
D\le D_n=\frac{L/2}{\tan\theta_n}.
\]
As franjas saem em pares simétricos ($\pm n$), por isso o número cai de $2$ em $2$: $7\to5\to3\to1$.

\textbf{a)} Como $y_n=D\tan\theta_n$ cresce com $D$, o número de franjas só diminui quando $D$ aumenta. Se o anteparo se aproximasse, as franjas convergiriam para o centro e o número não poderia cair. Como ele caiu de $7$ para $5$ e depois para $3$, \textbf{o anteparo está se afastando da rede}.

\textbf{b)} Em $t_1$ saem as ordens $\pm3$ ($7\to5$), ou seja, $y_3=L/2$:
\[
D(t_1)=\frac{0{,}6}{3/\sqrt7}=0{,}2\sqrt7=0{,}2\cdot 2{,}65\approx 0{,}53\ \text{m}.
\]

\textbf{c)} Em $t_1+2$ s saem as ordens $\pm2$ ($5\to3$):
\[
D(t_1+2)=0{,}6\sqrt3=0{,}6\cdot1{,}73\approx 1{,}038\ \text{m}.
\]
O movimento é retilíneo e uniforme, então
\[
v=\frac{\Delta D}{\Delta t}=\frac{0{,}6\sqrt3-0{,}2\sqrt7}{2}=\frac{1{,}038-0{,}530}{2}\approx 0{,}254\ \text{m/s}\approx 0{,}25\ \text{m/s}.
\]

\textbf{d)} Fica apenas o máximo central quando saem as ordens $\pm1$, ou seja, quando $D\tan\theta_1=L/2$:
\[
D_1=0{,}6\sqrt{15}=0{,}6\cdot3{,}87\approx 2{,}322\ \text{m}.
\]
A ordem $0$ ($\theta=0$) fica sempre no centro do anteparo e nunca sai. A distância percorrida desde $t_1$ é
\[
\Delta s=D_1-D(t_1)=0{,}6\sqrt{15}-0{,}2\sqrt7\approx 2{,}322-0{,}530\approx 1{,}79\ \text{m}.
\]
Isso corresponde a cerca de $1{,}792/0{,}254\approx 7{,}1$ s após $t_1$.

\textbf{Conferência.} As distâncias limite estão em ordem crescente, $D_3\approx0{,}53\ \text{m}<D_2\approx1{,}04\ \text{m}<D_1\approx2{,}32\ \text{m}$, como exige a sequência $7\to5\to3\to1$. As ordens de grandeza (décimos de m/s, metros) são compatíveis com uma montagem de laboratório.

ANSWER: a) Afastando-se: $y_n=D\tan\theta_n$ cresce com $D$, então as franjas saem pelas bordas; b) $D(t_1)=0{,}2\sqrt7\approx0{,}53$ m; c) $v=0{,}3\sqrt3-0{,}1\sqrt7\approx0{,}25$ m/s; d) $\Delta s=0{,}6\sqrt{15}-0{,}2\sqrt7\approx1{,}79$ m%
\newpage%
\section*{Questão 9}

\section*{Solução}

\textbf{Hipóteses.} O fóton se move em linha reta com velocidade $c$; na reflexão vale a lei da reflexão especular em relação à normal da esfera no ponto de impacto, tratada como se a superfície estivesse parada (como manda o enunciado).

\textbf{1) Direção da normal no ponto de impacto.} A reta $x=\sqrt{3}\,y$ faz $30^\circ$ com o eixo $x$. Os versores de incidência e de reflexão são
\[
\hat d=\left(\tfrac{\sqrt3}{2},\tfrac12\right),\qquad \hat d'=(0,-1).
\]
Na reflexão especular, $\hat d'=\hat d-2(\hat d\cdot\hat n)\hat n$, logo $\hat n\parallel \hat d'-\hat d=\left(-\tfrac{\sqrt3}{2},-\tfrac32\right)\parallel(1,\sqrt3)$. A normal \emph{externa} deve opor-se ao movimento do fóton que chega ($\hat d\cdot\hat n<0$), então
\[
\hat n=\left(-\tfrac12,-\tfrac{\sqrt3}{2}\right),\qquad \hat d\cdot\hat n=-\tfrac{\sqrt3}{2}=-\cos 30^\circ .
\]
Assim, o ângulo de incidência é $30^\circ$ (a normal faz $240^\circ$ com o eixo $x$, e é bissetriz entre $-\hat d$, a $210^\circ$, e $\hat d'$, a $270^\circ$).

\textbf{2) Condição de encontro.} No instante $t$ do impacto, o fóton está em $P=ct\,\hat d$, e esse ponto pertence à esfera na posição dada pela normal: $P=C+R\,\hat n$, com $C=(L,h)$ e $R=ut$. Por componentes:
\[
\frac{\sqrt3}{2}\,ct=L-\frac{R}{2},\qquad \frac12\,ct=h-\frac{\sqrt3}{2}\,R .
\]
Resolvendo o sistema linear em $ct$ e $R$:
\[
ct=\sqrt3\,L-h,\qquad R=ut=\sqrt3\,h-L .
\]
(Interpretação geométrica: a distância do centro à reta do fóton é $\tfrac{\sqrt3h-L}{2}=R\sin30^\circ$, e o percurso até o impacto vale a projeção $\hat d\cdot C=\tfrac{\sqrt3L+h}{2}$ menos $R\cos30^\circ$.)

\textbf{3) É de fato o primeiro contato.} Seja $g(t)=|ct\,\hat d-C|^2-(ut)^2$, um trinômio com coeficiente líder $c^2-u^2>0$ e $g(0)=L^2+h^2>0$ (fóton fora da esfera). No ponto achado, $g'(t)=2Rc(\hat d\cdot\hat n)-2uR=-2R\left(\tfrac{\sqrt3}{2}c+u\right)<0$: o fóton está \emph{entrando}, ou seja, trata-se da menor raiz de $g$, isto é, do primeiro encontro. Logo não há reflexão anterior.

\textbf{a)} Dividindo as duas relações:
\[
\boxed{\,u=c\,\frac{\sqrt3\,h-L}{\sqrt3\,L-h}\,}
\qquad\left(t=\frac{\sqrt3L-h}{c},\ R=\sqrt3h-L\right).
\]
Unidades: razão de comprimentos vezes $c$, ou seja, velocidade. Conferência numérica: $L=1$, $h=0{,}8$ dá $u\approx0{,}414\,c$, e a reflexão calculada sai exatamente na direção $(0,-1)$.

\textbf{b)} São necessários $t>0$, $R=ut>0$ (esfera que de fato se expande, $u>0$) e $u<c$ (dado).
\begin{itemize}
\item $t>0 \Rightarrow \sqrt3L-h>0 \Rightarrow h<\sqrt3L$;
\item $R>0 \Rightarrow \sqrt3h-L>0 \Rightarrow h>L/\sqrt3$. O centro precisa estar acima da reta do fóton, como na figura;
\item $u<c$: como o denominador é positivo, $\sqrt3h-L<\sqrt3L-h \Leftrightarrow (\sqrt3+1)h<(\sqrt3+1)L \Leftrightarrow h<L$.
\end{itemize}
Juntando (a condição $h<L$ é mais forte que $h<\sqrt3L$):
\[
\boxed{\,\frac{L}{\sqrt3}<h<L\,}\qquad\text{ou seja}\qquad \frac{\sqrt3}{3}<\frac hL<1 .
\]
Geometricamente, o centro deve ficar entre a reta do fóton ($y=x/\sqrt3$, a $30^\circ$) e a bissetriz $y=x$ ($45^\circ$).

\textbf{Casos-limite.} Se $h\to L/\sqrt3$, então $u\to0$: o centro fica sobre a trajetória e a esfera degenera num ponto. Se $h\to L$, então $u\to c$. A razão $h/L=\dfrac{c+\sqrt3u}{\sqrt3c+u}$ cresce monotonamente com $u$, então todo o intervalo é atingível.

ANSWER: a) $u=c\,\dfrac{\sqrt3\,h-L}{\sqrt3\,L-h}$; b) $\dfrac{L}{\sqrt3}<h<L$.%
\newpage%
\section*{Questão 10}

\section*{Solução}

\textbf{Referencial e hipóteses.} Tomo $x$ horizontal para a direita, com origem no eixo central, $y$ vertical para cima e $z$ saindo do papel. Os fios ficam em $x=\mp d/2$. Pela lei de Ampère, cada fio infinito gera $B=\mu_0 i/(2\pi r)$, perpendicular ao plano dos fios. Assim, no plano vale $\vec B=B_z\hat z$, e a força $\vec F=q\vec v\times\vec B$ também fica no plano, o que é coerente com a observação do enunciado. Pela regra da mão direita, à direita de uma corrente para cima o campo entra no papel; à esquerda, sai. Desprezo o peso, pois $g$ não é dado. A força magnética não realiza trabalho, logo $|\vec v|=v$ é constante. Com $\vec v=(v_x,v_y)$:
\[
\vec F=q\,(v_y B_z,\;-v_x B_z)\quad\Rightarrow\quad F_x=q\,v_yB_z .
\]

\textbf{a) Correntes opostas (situações 1 e 4).} Entre os fios, os dois campos entram no papel. Usando $1/(1+x)\approx 1-x$:
\[
B_z(x)=-\frac{\mu_0 i}{\pi d}\left[\frac{1}{1+2x/d}+\frac{1}{1-2x/d}\right]\approx-\frac{2\mu_0 i}{\pi d}\neq 0 .
\]
Já sobre o eixo atua $F_x=q\,v_yB_z\neq0$, que, em primeira ordem, não depende de $x$:
\begin{itemize}
\item situação 1 ($v_y=-v$): $F_x=+\dfrac{2\mu_0 q i v}{\pi d}$, para a direita;
\item situação 4 ($v_y=+v$): a mesma força, para a esquerda.
\end{itemize}
Logo o movimento vertical pelo eixo não é possível. A partícula descreve um arco de raio aproximado $R\approx \pi m v d/(2\mu_0 q i)$ e se desvia para um dos fios. Não existe termo restaurador, e as duas situações ficam descartadas.

\textbf{Correntes iguais (situações 2 e 3).} Os campos dos dois fios têm sentidos opostos:
\[
B_z(x)=\frac{\mu_0 i}{\pi d}\left[\frac{1}{1-2x/d}-\frac{1}{1+2x/d}\right]\approx\frac{\mu_0 i}{\pi d}\cdot\frac{4x}{d}=\frac{4\mu_0 i}{\pi d^2}\,x .
\]
Como $B(0)=0$, o movimento retilíneo sobre o eixo é possível. Para um pequeno desvio $x$:
\[
F_x=q\,v_y\,\frac{4\mu_0 i}{\pi d^2}\,x .
\]
\begin{itemize}
\item \emph{Situação 3} ($v_y=+v$): $F_x=+Kx$, que afasta a partícula do eixo, então o equilíbrio é instável. Fisicamente, a carga positiva subindo equivale a uma corrente paralela às dos fios e por isso é atraída por eles, com mais força pelo fio mais próximo. Situação descartada.
\item \emph{Situação 2} ($v_y=-v$): $F_x=-Kx$, com $K=\dfrac{4\mu_0 q i v}{\pi d^2}>0$. A força é restauradora. Fisicamente, a carga equivale a uma corrente antiparalela e é repelida pelos fios, com mais força pelo fio mais próximo.
\end{itemize}
\textbf{Somente a situação 2 satisfaz a condição.}

\textbf{b)} Suponho amplitude $A\ll d$, de modo que $v_x\ll v$. Então $v_y=-\sqrt{v^2-v_x^2}\approx -v$ e $F_y=-q v_xB_z$ são correções de segunda ordem. Resta
\[
m\ddot x=-Kx\;\Rightarrow\;\omega=\sqrt{\frac{K}{m}}=\frac{2}{d}\sqrt{\frac{\mu_0 q i v}{\pi m}} ,
\]
ou seja, um MHS em $x$ enquanto a partícula desce com velocidade $v$. A partícula cruza o eixo $y$ duas vezes por período. Na figura, $p$ corresponde a um lóbulo, isto é, a um intervalo de $T/2=\pi/\omega$ entre cruzamentos consecutivos. Assim,
\[
p=v\,\frac{T}{2}=\frac{\pi v}{\omega}=\frac{\pi d}{2}\sqrt{\frac{\pi m v}{\mu_0 q i}} .
\]

\textbf{Verificações.}
\begin{itemize}
\item Unidades: $[\mu_0 i\,q v]=(\mathrm{T\,m})(\mathrm{C})(\mathrm{m/s})=\mathrm{N\,m}$, logo $\mu_0 q i/(m v)$ é adimensional e $p$ sai em metros, proporcional a $d$.
\item Tendências: aumentar $i$ deixa a ``mola'' mais rígida e diminui $p$. Já $p\propto\sqrt{v}$, porque a rigidez cresce com $v$ mas o avanço vertical também cresce.
\item Para pequenas amplitudes, $p$ não depende da amplitude (isocronismo).
\end{itemize}

ANSWER: a) Apenas a situação 2. Nas situações 1 e 4, $B(0)\neq0$ e há força lateral já no eixo; na situação 3, a força afasta a partícula do eixo. b) $p\approx \dfrac{\pi d}{2}\sqrt{\dfrac{\pi m v}{\mu_0 q i}}$%
\newpage%
\end{document}